\documentclass[../../main.tex]{subfiles}
\begin{document}


{\bf [解]}
まず与えられた$f(t)=at+b$の条件から$a$と$b$を決定する．
\begin{align}
& \int_{-1}^1 f(t) \, dt = 2 \int_0^1 b \, dt = 2b = 1 \therefore b = \dfrac{1}{2} \\
& \int_{-1}^1 x f(x) \, dx = 2 \int_0^1 a x^2 \, dx = \dfrac{2}{3} a = 0 \therefore a = 0
\end{align}
だから, 
\begin{align}
 f(x) = \dfrac{1}{2} \label{eq:1}
\end{align}
である．これを$G,H$の式に代入すると
\begin{align}
 G(t) = \dfrac{1}{2} (t+1) \label{eq:2}
\end{align}
および
\begin{align}
 A = \int_{-1}^1 \dfrac{1}{2} x^2 \, dx = \int_0^1 x^2 \, dx = \dfrac{1}{3}
\end{align}
より
\begin{align}
 H(t) = \dfrac{1}{1 + 3t^2} \label{eq:3}
\end{align}
である．\cref{eq:2,eq:3}の大小を比べる．$-1\le t\le 0$では$G\ge 0, H>0$だから
\begin{align}
F(t) = \dfrac{G(t)}{H(t)} 
\end{align}
が1より大きいか小さいかで判別できる．
\begin{align}
 F(t) = \dfrac{1}{2}(t+1)(3t^2+1) = \dfrac{1}{2}(3t^3 + 3t^2 + t + 1) 
\end{align}
より一階微分は
\begin{align}
 F'(t) = \dfrac{(3t+1)^2}{2} \geqq 0
\end{align}
だから$F(t)$は区間内で単調増加で，$F(0)=\dfrac{1}{2}$から$-1 \leqq t \leqq 0$ では
\begin{align}
 & F(t) \leqq \dfrac{1}{2} < 1 \\
\therefore 
 & G(t) < H(t)
\end{align}
である．

\end{document}
